\chapter{A Continuum-Mechanics Primer}
\label{ch:continuum}

The finite element method solves the field equations of continuum
mechanics. This chapter states those equations for the linear
(small-strain, linear-elastic) theory that underlies the great majority
of engineering stress analysis, and casts them in the matrix (Voigt)
form the FEM will consume. Readers fluent in elasticity may skim to
\cref{ch:strong-weak}. The development follows the standard
treatments~\citep{malvern1969,hughes2000,bathe2014}.

\section{Kinematics: the strain tensor}

Let $\disp(\vec x)$ be the displacement field of a body occupying the
domain $\domain\subset\mathbb R^3$. Under the \emph{small-strain}
assumption $\lVert\grad\disp\rVert\ll 1$, deformation is measured by the
infinitesimal (Cauchy) strain tensor, the symmetric part of the
displacement gradient:
\begin{equation}
  \varepsilon_{ij} = \tfrac12\!\left(\pd{u_i}{x_j}+\pd{u_j}{x_i}\right),
  \qquad
  \strain = \tfrac12\bigl(\grad\disp+\grad\disp\transpose\bigr)
          = \sym\grad\disp .
  \label{eq:strain-disp}
\end{equation}
Equation~\eqref{eq:strain-disp} is the \emph{strain--displacement
relation}; it is linear in $\disp$, which is precisely what makes the
resulting problem linear. Discarding the antisymmetric part (the
infinitesimal rotation) is legitimate only when strains \emph{and}
rotations are small; large-rotation problems require the nonlinear
strain measures of \cref{ch:extensions}.

\section{Kinetics: stress and equilibrium}

Internal forces are described by the \textbf{Cauchy stress tensor}
$\stress$. By Cauchy's theorem the traction on an internal surface with
unit normal $\vec n$ is
\begin{equation}
  \traction = \stress\,\vec n, \qquad t_i = \sigma_{ij}n_j .
\end{equation}
Balance of angular momentum makes $\stress$ symmetric
($\sigma_{ij}=\sigma_{ji}$), leaving six independent components. Balance
of linear momentum, in the static case, gives the \textbf{equilibrium
equations}
\begin{equation}
  \pd{\sigma_{ij}}{x_j} + b_i = 0, \qquad
  \divergence\stress + \bodyforce = \vec 0 \quad\text{in }\domain,
  \label{eq:equilibrium}
\end{equation}
where $\bodyforce$ is the body force per unit volume. (The dynamic case
replaces the right-hand side with $\rho\,\ddot{\disp}$; see
\cref{ch:extensions}.)

\section{Constitutive law: linear isotropic elasticity}

For a linear elastic material, stress is linear in strain via the
fourth-order elasticity tensor, $\sigma_{ij}=C_{ijkl}\varepsilon_{kl}$.
For the \emph{isotropic} case this reduces to two constants. In terms of
the Lam\'e parameters $\lambda,\mu$,
\begin{equation}
  \sigma_{ij} = \lambda\,\varepsilon_{kk}\,\delta_{ij} + 2\mu\,\varepsilon_{ij},
  \label{eq:hooke}
\end{equation}
with $\mu=G=E/[2(1+\nu)]$ the shear modulus and
$\lambda=E\nu/[(1+\nu)(1-2\nu)]$, where $E$ is Young's modulus and $\nu$
Poisson's ratio. The inverse relation is
$\varepsilon_{ij}=\tfrac{1}{E}[(1+\nu)\sigma_{ij}-\nu\sigma_{kk}\delta_{ij}]$.

\section{Voigt form and the elasticity matrix}

For computation we collapse the symmetric tensors to six-vectors
(\cref{app:notation}) and write $\svec=\Dmat\,\evec$ with the strain
vector using \emph{engineering} shear strains $\gamma_{ij}=2\varepsilon_{ij}$.
The 3D isotropic elasticity matrix is
\begin{equation}
  \Dmat = \frac{E}{(1+\nu)(1-2\nu)}
  \begin{bmatrix}
    1-\nu & \nu & \nu & 0 & 0 & 0\\
    \nu & 1-\nu & \nu & 0 & 0 & 0\\
    \nu & \nu & 1-\nu & 0 & 0 & 0\\
    0 & 0 & 0 & \tfrac{1-2\nu}{2} & 0 & 0\\
    0 & 0 & 0 & 0 & \tfrac{1-2\nu}{2} & 0\\
    0 & 0 & 0 & 0 & 0 & \tfrac{1-2\nu}{2}
  \end{bmatrix}.
  \label{eq:Dmat3d}
\end{equation}
The engineering-shear convention is what makes $\svec\transpose\evec$ the
correct (scalar) work-conjugate product used to form the strain energy.

\subsection{Two-dimensional reductions}

For thin bodies loaded in-plane (\textbf{plane stress},
$\sigma_{zz}=\tau_{xz}=\tau_{yz}=0$),
\begin{equation}
  \Dmat = \frac{E}{1-\nu^2}
  \begin{bmatrix} 1 & \nu & 0\\ \nu & 1 & 0\\ 0 & 0 & \tfrac{1-\nu}{2}\end{bmatrix},
  \qquad \evec=[\varepsilon_{xx},\varepsilon_{yy},\gamma_{xy}]\transpose .
\end{equation}
For long prismatic bodies (\textbf{plane strain},
$\varepsilon_{zz}=\gamma_{xz}=\gamma_{yz}=0$),
\begin{equation}
  \Dmat = \frac{E}{(1+\nu)(1-2\nu)}
  \begin{bmatrix} 1-\nu & \nu & 0\\ \nu & 1-\nu & 0\\ 0 & 0 & \tfrac{1-2\nu}{2}\end{bmatrix},
\end{equation}
with the out-of-plane stress recovered as
$\sigma_{zz}=\nu(\sigma_{xx}+\sigma_{yy})\neq 0$.

\section{Thermal strains}

A temperature change $\Delta T$ adds an isotropic \emph{initial}
(thermal) strain $\evec_0=\alpha\,\Delta T\,[1,1,1,0,0,0]\transpose$,
where $\alpha$ is the coefficient of thermal expansion. Stress is then
driven only by the \emph{mechanical} (total minus thermal) strain:
\begin{equation}
  \svec = \Dmat\,(\evec-\evec_0).
  \label{eq:thermal-strain}
\end{equation}
This term reappears in \cref{ch:elements-assembly} as a thermal load
vector, and it is the mechanism behind thermomechanical (coupled)
analyses in \cref{ch:extensions}.

\section{The boundary-value problem, informally}

Collecting \eqref{eq:strain-disp}, \eqref{eq:equilibrium} and
\eqref{eq:hooke}, elastostatics asks: find $\disp$ on $\domain$
satisfying equilibrium, with the material law and strain definition
substituted in, subject to prescribed displacements on part of the
boundary and prescribed tractions on the rest. Substituting throughout
yields the \textbf{Navier--Lam\'e} displacement equations
\begin{equation}
  \mu\,\grad^2\disp + (\lambda+\mu)\,\grad(\divergence\disp) + \bodyforce = \vec 0 .
  \label{eq:navier}
\end{equation}
The next chapter states this boundary-value problem precisely (its
\emph{strong form}) and then recasts it into the \emph{weak form} that
the finite element method actually discretises.
